\documentclass[11pt,a4paper]{article}
\usepackage[margin=2.1cm]{geometry}
\usepackage{parskip}
\usepackage{xcolor}
\usepackage{titlesec}
\usepackage{enumitem}
\usepackage{fancyhdr}
\usepackage[colorlinks=true,linkcolor=black,urlcolor=black]{hyperref}
\usepackage[breakable]{tcolorbox}

\definecolor{ink}{HTML}{1C1C1A}
\definecolor{accent}{HTML}{2A78D6}
\definecolor{soft}{HTML}{5A5A55}
\definecolor{rule}{HTML}{D5D5CF}

\color{ink}
\titleformat{\section}{\large\bfseries\color{accent}}{\thesection.}{0.5em}{}
\titleformat{\subsection}{\normalsize\bfseries}{\thesubsection}{0.5em}{}
\titlespacing{\section}{0pt}{14pt}{6pt}
\setlist[itemize]{leftmargin=1.2em,itemsep=2pt,topsep=3pt}

\pagestyle{fancy}\fancyhf{}
\fancyhead[L]{\footnotesize\color{soft}Research-loop prompt for a new ML project}
\fancyhead[R]{\footnotesize\color{soft}\thepage}
\renewcommand{\headrulewidth}{0.4pt}

\newtcolorbox{promptbox}{colback=white,colframe=rule,boxrule=0.8pt,
  left=9pt,right=9pt,top=9pt,bottom=9pt,breakable,
  before upper={\setlength{\parskip}{5pt}}}

\begin{document}

\begin{center}
{\LARGE\bfseries A prompt for running a machine-learning project}\\[3pt]
{\large\bfseries with GitHub, Kaggle GPU and a self-correcting research loop}\\[8pt]
{\color{soft}\small Distilled from a completed project: gait-based Parkinson's
classification, 20 GPU sessions, 48 configurations.\\
Written so the next project repeats the method and not the mistakes.}
\end{center}

\vspace{4pt}
\hrule
\vspace{10pt}

\section{How to use this}

Paste Part~A as your first message in a new session, after filling in the four
placeholders. Parts~B and~C are the standing rules; paste them too, or save them
as \texttt{CLAUDE.md} in the project folder so they load automatically every
session. Part~D is what to say when you want the loop to continue.

\section{Part A --- the opening message}

\begin{promptbox}
\small
I am starting a new machine-learning project and I want you to run it as a
senior researcher would, not as a code monkey. Read this whole message before
touching anything.

\vspace{3pt}\textbf{THE PROBLEM}

$\bullet$ Goal: \texttt{<what must be predicted, and why it matters>}\\
$\bullet$ Data: \texttt{<what I have, how much, where it lives, its size>}\\
$\bullet$ Constraint: \texttt{<required model family, deadline, or "none">}\\
$\bullet$ Success looks like: \texttt{<the decision this model should support>}

\vspace{3pt}\textbf{THE INFRASTRUCTURE I WANT}

Set this up first, before any modelling:

1. A GitHub repository for the code. You commit and push after every meaningful
change, with messages that say why, not what. Data and outputs stay out of git.

2. A Kaggle notebook that pulls the latest commit and trains on the GPU. I want
to reproduce any result by pressing Run All. Put the data in a private Kaggle
dataset; keep the raw data out of the repository.

3. Every run writes a structured \texttt{results.json} --- metrics, the exact
configuration, environment versions, runtime --- plus its full log. Both get
archived under version control of some kind, so a run from three weeks ago is
still readable.

4. After each run you fetch the log and the JSON back, parse them, and tell me
what happened in plain language: what improved, what did not, and what the
numbers actually license you to claim. Then you decide the next experiment from
that evidence and push the next iteration.

Ask me for the Kaggle token when you need it; I will paste a fresh one. Never
commit it.

\vspace{3pt}\textbf{HOW I WANT YOU TO WORK}

Plan like someone with a doctorate in machine learning. That means: state a
hypothesis before each experiment, say what result would falsify it, run the
cheapest test that could falsify it, and update. Do not run a grid because a
grid is available.

Before spending GPU time on an idea, look for a five-minute version of it that
runs locally. A quick probe on aggregated features usually predicts whether a
representation carries signal, and it costs nothing.

When you stall --- roughly three consecutive experiments with no improvement
that survives a significance test --- stop tuning and go read. Search for recent
peer-reviewed work on this exact problem and on the methods you are using.
Bring me a short summary of what those papers do differently, and a concrete
proposal for what to adopt. If a paper is paywalled, give me the title, authors,
year and DOI and I will fetch the PDF for you.

Tell me plainly when you hit a ceiling, and diagnose which kind it is: not
enough data, not enough signal in the features, a flaw in how the data was
collected or split, or a genuine limit of the method. Each has a different
remedy and only one of them is "train longer".

\vspace{3pt}\textbf{WHAT I DO NOT WANT}

Do not report a number you have not computed. Do not present a result whose
difference from the baseline is inside the noise as an improvement. Do not
quietly change the evaluation protocol to make a number look better. If an
experiment fails, tell me it failed and what you learned; a negative result that
saves me a week is worth more than an optimistic one that costs me a month.

Start by proposing the evaluation protocol and the infrastructure. Do not train
anything until I have agreed to the protocol.
\end{promptbox}

\section{Part B --- the standing scientific rules}

\begin{promptbox}
\small
These rules hold for the whole project. They exist because each one has already
cost someone a wasted month.

\vspace{3pt}\textbf{FIX THE EVALUATION PROTOCOL BEFORE THE FIRST EXPERIMENT}

Decide the unit of independence --- the patient, the subject, the session, the
device, the site --- and split on it. Never split on rows, windows or augmented
copies. Two windows cut a second apart from one recording are near duplicates:
put one in train and one in test and the model will recognise the individual
instead of the phenomenon, and the score will look wonderful.

Quantify this once, early, and keep the number. Re-run the same pipeline with a
naive random-row split and report both. In my last project the same Random
Forest scored 73\,\% under subject-wise splits and 98\,\% under row-wise splits.
That single comparison explained a 30-point spread in the published literature.

Write the split assignment to a file and load it from there, so every model in
the project sees identical folds.

\vspace{3pt}\textbf{FIT NOTHING ON DATA YOU WILL EVALUATE ON}

Normalisation statistics, feature selection, imputation values, class weights,
decision thresholds: all estimated on training data only, then applied
unchanged. At deployment, ship the training statistics alongside the weights and
apply those --- never statistics recomputed from the new cohort.

\vspace{3pt}\textbf{SELECT ON VALIDATION, NEVER ON TEST}

This includes the choice between experiments, not just hyper-parameters inside
one run. If you compare forty configurations by their test score and report the
best, the test set has entered the model definition and the number is
optimistic. Rank configurations by validation performance and report the winner
of that ranking.

Measure the cost of getting this wrong and report it. In my last project,
choosing on test gave 78\,\%, choosing on validation gave 75\,\%; those three
points were pure optimism.

\vspace{3pt}\textbf{ALWAYS BUILD A TRIVIAL BASELINE FIRST}

Majority class, one hand-computed feature, a linear model on aggregates. If the
deep model does not clearly beat it, that is the headline finding and you should
say so. In my last project a single duty-cycle scalar matched a tuned
state-space model to within noise --- which told us the signal was saturated,
not that the architecture was wrong. We would have wasted weeks without it.

\vspace{3pt}\textbf{REPORT UNCERTAINTY, AND RESPECT IT}

Every number gets a bootstrap interval. If the rows are correlated --- several
windows per subject --- resample whole groups, not rows; resampling rows treats
duplicates as independent evidence and halves the apparent uncertainty.

To compare two models, use a paired bootstrap over the same units and report the
interval of the difference. If it contains zero, say the models are
indistinguishable. Do not rank them anyway.

Know the resolution of your experiment. With one hundred subjects the 95\,\%
interval on an accuracy near 0.75 is roughly $\pm 8$ points, so a
three-point difference is unmeasurable no matter how many seeds you average.

\vspace{3pt}\textbf{VERIFY, DO NOT ASSUME}

Check invariants empirically rather than by reading code. Reconstruct a
normalised test split by hand from training statistics and compare it against
what the pipeline produced. Confirm the same split is \emph{not} reproducible
from all-data statistics. Recompute one sample's features in isolation and check
they match the batch-built cache byte for byte. Assert that no unit appears in
two splits. Keep these as a script and run it before every submission.

\vspace{3pt}\textbf{REPRODUCE ANY NUMBER BEFORE ACTING ON IT}

This applies to numbers from papers, from me, and from any expert or reviewer
you consult. In my last project a reviewing agent reported that fusing two
models reached 0.898 AUC; reproducing it gave 0.856 against a 0.854 baseline ---
nothing. The same review claimed a missing feature family was the key lever; a
five-minute probe showed it added nothing. Both would have cost GPU hours and a
wrong conclusion in the report.
\end{promptbox}

\section{Part C --- the loop, and what to do when it stalls}

\begin{promptbox}
\small
\vspace{3pt}\textbf{THE ITERATION}

Each cycle is: hypothesis $\rightarrow$ cheapest falsifying test $\rightarrow$
run $\rightarrow$ fetch log $\rightarrow$ analyse $\rightarrow$ decide. Report to
me after each cycle in a few sentences: what you expected, what happened,
whether the difference is significant, and what you will try next.

Keep a running knowledge-base file in the repository --- call it
\texttt{CLAUDE.md} --- and update it as you go, not at the end. It should hold:
the current result, the protocol, every configuration tried with its numbers,
every bug hit and its fix, the timings, and a section named "things not to redo".
Write it for a stranger, or for me in a year. When a session ends, that file is
the only thing that survives.

Archive every run's \texttt{results.json} and log permanently, outside temporary
directories. You will want to compare against run 7 when you are on run 30.

\vspace{3pt}\textbf{WHEN THREE EXPERIMENTS IN A ROW FAIL TO IMPROVE ANYTHING}

Stop tuning. In order:

1. Re-examine the data, not the model. Plot the classes. Compute what the
strongest single feature achieves. Check whether some nuisance variable ---
site, device, study, acquisition date, sequence length --- is correlated with
the label and with the split. In my last project the three sub-studies had
different class rates and the prescribed folds were not stratified by study, so
"which study" was a shortcut that helped in training and scored below chance at
test. That explained the fold-to-fold collapses that no optimiser setting fixed.

2. Go to the literature. Find recent peer-reviewed work on this problem and on
the method. Read for what they do differently in the parts you have been
ignoring: preprocessing, representation, the evaluation protocol, the cohort
size. Summarise for me: what they claim, under what protocol, and whether the
comparison to our setup is fair. Papers reporting far higher numbers very often
use a weaker split --- check before you believe them.

3. Propose the two or three most promising adaptations, with a cheap test for
each and a prediction of what it should produce. If a paper is behind a
paywall, send me the title, authors, year and DOI and I will get it for you.

\vspace{3pt}\textbf{DIAGNOSING THE CEILING}

When you believe you have reached the limit, say so plainly and support it. The
diagnosis matters more than the number, because each cause has a different
remedy:

$\bullet$ \emph{Sample size}: several very different model families land in the
same narrow band, and no difference between them is significant. Remedy: more
units, not more capacity. Estimate how many you would need for the effect you
want to detect, and tell me.

$\bullet$ \emph{Signal saturation}: a trivial feature reaches nearly the same
score as the deep model. Remedy: better measurements or richer features, not a
bigger model.

$\bullet$ \emph{Protocol or cohort flaw}: performance varies wildly across folds
in a way tied to some grouping. Remedy: fix the split or acknowledge the
confound; often it cannot be fixed and must be reported.

$\bullet$ \emph{Optimisation}: training diverges or collapses, best epoch is
epoch one, loss curves look pathological. This is the only case where the
remedy is in the training loop.

Tell me which one it is, what evidence supports it, and what the honest next
research step would be. If the answer is "we need more data", say it --- and
tell me what data, how much, and where it might come from.

\vspace{3pt}\textbf{DELIVERABLES}

Keep the write-up generated from the run outputs, never typed by hand. One
script should regenerate every table and figure from the archived JSON, so the
document cannot drift from the experiments. Numbers appearing in two places must
be produced by the same code path, and it is worth writing a check that compares
them.
\end{promptbox}

\section{Part D --- what to say to continue the loop}

\begin{promptbox}
\small
Continue the loop. Fetch the last run's log and results, analyse them, tell me
whether the change was significant under a paired bootstrap, and decide the next
experiment from the evidence. If we have now had three cycles without a
significant improvement, switch to the literature step and come back with
proposals rather than another run.
\end{promptbox}

\section{Part E --- the four placeholders in Part A}

\begin{itemize}
\item \textbf{Goal} --- what is predicted and what decision it supports. "Detect
X from Y so that a clinician can prioritise a scan" is useful; "get high
accuracy" is not.
\item \textbf{Data} --- how many independent units (patients, subjects, sites),
how much per unit, the format, and where the files are. State the natural unit
of independence explicitly; it determines the split.
\item \textbf{Constraint} --- a required model family, a deadline, a compute
budget, a report template. Write "none" if there is none.
\item \textbf{Success} --- the target if you have one, and whether a modest
honest number beats an impressive fragile one. Say which, because it changes
every decision that follows.
\end{itemize}

\vspace{6pt}
\hrule
\vspace{6pt}
{\footnotesize\color{soft}
The lessons in Parts B and C are not hypothetical: each corresponds to something
that went wrong, was caught, and was fixed in a real project. The most expensive
were selecting configurations on the test set, an uncertainty interval that was
twice too narrow because correlated rows were resampled independently, and two
plausible suggestions from an expert reviewer that evaporated on reproduction.}

\end{document}
